\chapter{Stabilität}

\section{Die Frage}

In den beiden vorigen Kapiteln haben wir Gleichgewichtspunkte gefunden und ihr Verhalten aus Bildern abgelesen. Das genügt bei zwei Gleichungen mit geraden Nullklinen. Bei allem Weiteren nicht mehr.

\begin{merksatz}
\boxtitel{Was wir brauchen}
\begin{center}
\begin{minipage}{0.9\textwidth}
\textbf{Ein Rechenverfahren, das aus den Gleichungen bestimmt, ob ein Gleichgewicht stabil ist --- ohne Bild und ohne Lösung.}
\end{minipage}
\end{center}

\schritt{Der Gedanke}
Er steht bereits in Satz 3.4: In der Nähe eines Gleichgewichts ersetzt man die Funktion durch ihre Tangente und schaut nach, ob kleine Abweichungen wachsen oder schrumpfen.

Neu ist nur, dass wir zwei Größen gleichzeitig betrachten --- und dass \enquote{Tangente} dann etwas anderes heißt.
\end{merksatz}

\section{Die Jacobi-Matrix}

\begin{definition}
\boxtitelm{Definition 6.1}
Für das System
\[ \frac{\dd x}{\dd t} = f(x,y), \qquad \frac{\dd y}{\dd t} = g(x,y) \]
heißt
\[ J = \begin{pmatrix}
\dfrac{\partial f}{\partial x} & \dfrac{\partial f}{\partial y}\\[8pt]
\dfrac{\partial g}{\partial x} & \dfrac{\partial g}{\partial y}
\end{pmatrix} \]
die \Begriff{Jacobi-Matrix}. Für die Stabilitätsuntersuchung wird sie im Gleichgewichtspunkt ausgewertet.
\end{definition}

\begin{merksatz}
\boxtitel{Warum eine Matrix}
Sei $(x^{*}, y^{*})$ ein Gleichgewicht und
\[ x = x^{*} + u, \qquad y = y^{*} + v \]
mit kleinen Abweichungen $u, v$.

\schritt{Die Näherung}
Nach der Taylor-Entwicklung in zwei Veränderlichen (Band 3, Teil~V):
\[ f(x^{*}+u,\ y^{*}+v) \approx \underbrace{f(x^{*},y^{*})}_{= 0} + \frac{\partial f}{\partial x}\,u + \frac{\partial f}{\partial y}\,v \]

Der erste Term verschwindet, weil es ein Gleichgewicht ist. Also
\[ \frac{\dd u}{\dd t} \approx \frac{\partial f}{\partial x}\,u + \frac{\partial f}{\partial y}\,v , \qquad \frac{\dd v}{\dd t} \approx \frac{\partial g}{\partial x}\,u + \frac{\partial g}{\partial y}\,v \]

oder kurz
\[ \frac{\dd}{\dd t}\begin{pmatrix}u\\v\end{pmatrix} = J\begin{pmatrix}u\\v\end{pmatrix} . \]

\begin{center}
\begin{minipage}{0.88\textwidth}
\kursiv{Das ist genau die Gleichung aus Kapitel~1 --- $\dd N/\dd t = rN$ ---, nur mit einer Matrix statt einer Zahl. Und die Lösung sieht entsprechend aus: Sie besteht aus Exponentialfunktionen, deren Raten die \kursiv{Eigenwerte} der Matrix sind.}
\end{minipage}
\end{center}
\end{merksatz}

\begin{merksatz}
\boxtitel{Satz 6.2 --- das Stabilitätskriterium}
Seien $\lambda_1, \lambda_2$ die Eigenwerte von $J$ im Gleichgewichtspunkt. Dann gilt:

\begin{center}
{\small
\renewcommand{\arraystretch}{1.55}
\begin{tabular}{@{}P{3.4cm}P{5.0cm}@{}}
\toprule
\textbf{Beide Realteile $<0$} & \textcolor{gruen}{stabil} --- alle nahen Bahnen laufen hin\\
\textbf{Ein Realteil $>0$} & \textcolor{akzent}{instabil} --- es gibt eine Fluchtrichtung\\
\textbf{Ein Realteil $=0$} & \textcolor{mittel}{unentschieden} --- die Näherung reicht nicht\\
\bottomrule
\end{tabular}}
\end{center}

\schritt{Warum}
Die Lösung der linearisierten Gleichung hat die Gestalt
\[ u(t) = c_1e^{\lambda_1 t} + c_2e^{\lambda_2 t} . \]
Ein Term wächst genau dann, wenn der Realteil des Exponenten positiv ist.

\begin{center}
\begin{minipage}{0.88\textwidth}
\kursiv{Die dritte Zeile ist keine Bequemlichkeit, sondern eine echte Grenze: Bei Realteil null entscheiden die vernachlässigten Terme höherer Ordnung. Genau dieser Fall trat in Kapitel~4 auf --- und deshalb war das Lotka-Volterra-Modell dort so empfindlich.}
\end{minipage}
\end{center}
\end{merksatz}

\section{Spur und Determinante}

\begin{merksatz}
\boxtitel{Man braucht die Eigenwerte nicht auszurechnen}
Für eine $2\times2$-Matrix $J = \begin{pmatrix} p & q\\ r & s\end{pmatrix}$ setzt man
\[ T = \operatorname{spur} J = p + s, \qquad D = \det J = ps - qr . \]

\schritt{Die charakteristische Gleichung}
Die Eigenwerte sind die Lösungen von
\[ \lambda^2 - T\lambda + D = 0 , \qquad\text{also}\qquad \lambda_{1,2} = \frac{T \pm \sqrt{T^2-4D}}{2} . \]

\schritt{Die beiden Beziehungen}
\[ \lambda_1 + \lambda_2 = T, \qquad \lambda_1\cdot\lambda_2 = D \]
Das folgt unmittelbar aus dem Satz von Vieta.

\begin{center}
\begin{minipage}{0.88\textwidth}
\kursiv{Daraus lässt sich die Stabilität ablesen, ohne die Wurzel zu ziehen --- was bei Handrechnungen und bei Formeln mit Buchstaben einen erheblichen Unterschied macht.}
\end{minipage}
\end{center}
\end{merksatz}

\begin{merksatz}
\boxtitel{Satz 6.3 --- das Kriterium in $T$ und $D$}
\[ \text{stabil} \quad\Longleftrightarrow\quad T < 0 \ \text{ und } \ D > 0 \]

\schritt{Beweis --- Fall 1: reelle Eigenwerte ($T^2 \ge 4D$)}
\kursiv{$D > 0$} bedeutet $\lambda_1\lambda_2 > 0$: beide Eigenwerte haben dasselbe Vorzeichen.\\
\kursiv{$T < 0$} bedeutet $\lambda_1+\lambda_2 < 0$: dieses gemeinsame Vorzeichen ist negativ.

Umgekehrt: Sind beide negativ, so ist ihr Produkt positiv und ihre Summe negativ.

\schritt{Beweis --- Fall 2: komplexe Eigenwerte ($T^2 < 4D$)}
Dann ist
\[ \lambda_{1,2} = \frac{T}{2} \pm \mathrm{i}\,\frac{\sqrt{4D-T^2}}{2} , \]
also ist der Realteil beider Eigenwerte gleich $T/2$. Er ist genau dann negativ, wenn $T < 0$.

Die Bedingung $D > 0$ ist hier automatisch erfüllt: Aus $T^2 < 4D$ folgt $D > T^2/4 \ge 0$. $\blacksquare$

\schritt{Die Kurzfassung}
\begin{center}
\begin{minipage}{0.9\textwidth}
\textbf{$D$ entscheidet, ob es ein Sattel ist. $T$ entscheidet, ob es stabil ist. Die Diskriminante $T^2-4D$ entscheidet, ob es schwingt.}
\end{minipage}
\end{center}
\end{merksatz}

\begin{merksatz}
\boxtitel{Die Landkarte}
\begin{center}
\begin{tikzpicture}[scale=1.0]
\draw[grau!60, line width=0.7pt] (-3.4,0) -- (3.4,0);
\draw[grau!60, line width=0.7pt] (0,-1.6) -- (0,3.4);
\draw[akzent, line width=1.1pt, smooth, samples=60, domain=-3.3:3.3]
  plot (\x, {0.3*\x*\x});
\node[font=\scriptsize, color=akzent, anchor=west] at (2.4,3.1) {$T^2 = 4D$};
\node[font=\scriptsize, color=grau] at (3.2,-0.3) {$T$};
\node[font=\scriptsize, color=grau] at (-0.3,3.2) {$D$};
\node[font=\scriptsize, color=gruen, align=center] at (-1.9,2.6) {stabiler\\Strudel};
\node[font=\scriptsize, color=akzent, align=center] at (1.9,2.6) {instabiler\\Strudel};
\node[font=\scriptsize, color=gruen, align=center] at (-2.6,0.65) {stabiler\\Knoten};
\node[font=\scriptsize, color=akzent, align=center] at (2.6,0.65) {instabiler\\Knoten};
\node[font=\scriptsize, color=akzent] at (0,-0.9) {Sattel (stets instabil)};
\draw[tief, line width=1.3pt] (0,0.05) -- (0,3.3);
\node[font=\scriptsize, color=tief, anchor=west] at (0.12,1.55) {Zentrum};
\end{tikzpicture}
\end{center}

\begin{center}
{\small
\renewcommand{\arraystretch}{1.5}
\begin{tabular}{@{}P{2.4cm}P{2.6cm}P{3.0cm}@{}}
\toprule
\textbf{Typ} & \textbf{Bedingung} & \textbf{Verhalten}\\
\midrule
Sattel & $D < 0$ & \textcolor{akzent}{zwei Richtungen, eine flieht}\\
Knoten & $D>0$, $T^2>4D$ & direkt hin oder weg\\
Strudel & $D>0$, $T^2<4D$ & \textcolor{mittel}{spiralig, schwingend}\\
Zentrum & $D>0$, $T=0$ & \textcolor{tief}{geschlossene Bahnen}\\
\bottomrule
\end{tabular}}
\end{center}
\end{merksatz}

\section{Die beiden Modelle geprüft}

\begin{merksatz}
\boxtitel{Lotka-Volterra --- warum ein Zentrum}
\[ f = x(a-by), \qquad g = y(-c+dx) \]

\schritt{Die Ableitungen}
\begin{align*}
\frac{\partial f}{\partial x} &= a - by & \frac{\partial f}{\partial y} &= -bx\\
\frac{\partial g}{\partial x} &= dy & \frac{\partial g}{\partial y} &= -c + dx
\end{align*}

\schritt{Im Gleichgewichtspunkt $\left(\frac{c}{d},\ \frac{a}{b}\right)$}
\[ \frac{\partial f}{\partial x} = a - b\cdot\frac{a}{b} = 0, \qquad \frac{\partial g}{\partial y} = -c + d\cdot\frac{c}{d} = 0 \]
\[ \frac{\partial f}{\partial y} = -\frac{bc}{d}, \qquad \frac{\partial g}{\partial x} = \frac{da}{b} \]

\[ J = \begin{pmatrix} 0 & -\frac{bc}{d}\\[4pt] \frac{da}{b} & 0 \end{pmatrix} \]

\schritt{Spur und Determinante}
\[ T = 0 + 0 = 0 \]
\[ D = 0\cdot0 - \left(-\frac{bc}{d}\right)\left(\frac{da}{b}\right) = +\,\frac{bc}{d}\cdot\frac{da}{b} = a\,c \]

\schritt{Das Ergebnis}
$T = 0$ und $D = ac > 0$: \textbf{ein Zentrum.}
\[ \lambda_{1,2} = \pm\mathrm{i}\sqrt{ac} \]

\begin{center}
\begin{minipage}{0.88\textwidth}
\kursiv{Damit ist die Periode aus Kapitel~4 begründet: Bei rein imaginären Eigenwerten $\pm\mathrm{i}\omega$ schwingt das System mit der Kreisfrequenz $\omega = \sqrt{ac}$, also der Periode $2\pi/\sqrt{ac}$. \checkmark}
\end{minipage}
\end{center}

\schritt{Und warum das so empfindlich ist}
Der Fall $T = 0$ ist genau der unentschiedene Fall aus Satz 6.2. Die Linearisierung sagt \enquote{Zentrum}, aber sie kann es nicht garantieren --- die vernachlässigten Terme entscheiden.

\begin{center}
\begin{minipage}{0.9\textwidth}
\textbf{Im Bild von Abschnitt 6.3 liegt das Zentrum genau auf der senkrechten Achse.} Jede noch so kleine Änderung schiebt den Punkt nach links oder rechts --- und macht daraus einen Strudel.
\end{minipage}
\end{center}

\begin{center}
\begin{minipage}{0.88\textwidth}
\kursiv{Das ist die strukturelle Instabilität aus Kapitel~4, nun in einem Satz begründet: Zentren sind Grenzfälle, und Grenzfälle sind fragil. Beim Lotka-Volterra-Modell rettet die Erhaltungsgröße aus Satz 4.3 das Zentrum --- aber nur, solange man das Modell nicht anrührt.}
\end{minipage}
\end{center}
\end{merksatz}

\begin{merksatz}
\boxtitel{Das Konkurrenzmodell --- Satz 5.3 nachgewiesen}
\[ f = r_1x\left(1 - \frac{x+\alpha_{12}y}{K_1}\right) \]

\schritt{Die Ableitungen}
Ausmultipliziert ist $f = r_1x - \frac{r_1}{K_1}x^2 - \frac{r_1\alpha_{12}}{K_1}xy$, also
\[ \frac{\partial f}{\partial x} = r_1 - \frac{2r_1}{K_1}x - \frac{r_1\alpha_{12}}{K_1}y, \qquad \frac{\partial f}{\partial y} = -\frac{r_1\alpha_{12}}{K_1}x \]

\schritt{Im Koexistenzpunkt}
Dort gilt $x^{*}+\alpha_{12}y^{*} = K_1$. Wir schreiben die erste Ableitung um:
\[ \frac{\partial f}{\partial x} = r_1\left(1 - \frac{x^{*}+\alpha_{12}y^{*}}{K_1} - \frac{x^{*}}{K_1}\right) = r_1\left(1 - 1 - \frac{x^{*}}{K_1}\right) = -\frac{r_1x^{*}}{K_1} \]

Entsprechend für $g$. Insgesamt
\[ J = \begin{pmatrix}
-\frac{r_1x^{*}}{K_1} & -\frac{r_1\alpha_{12}x^{*}}{K_1}\\[6pt]
-\frac{r_2\alpha_{21}y^{*}}{K_2} & -\frac{r_2y^{*}}{K_2}
\end{pmatrix} \]

\schritt{Die Spur}
\[ T = -\frac{r_1x^{*}}{K_1} - \frac{r_2y^{*}}{K_2} \;<\; 0 \]
--- stets negativ, da alle Größen positiv sind.

\schritt{Die Determinante}
\begin{align*}
D &= \frac{r_1x^{*}}{K_1}\cdot\frac{r_2y^{*}}{K_2} - \frac{r_1\alpha_{12}x^{*}}{K_1}\cdot\frac{r_2\alpha_{21}y^{*}}{K_2}\\
&= \frac{r_1r_2x^{*}y^{*}}{K_1K_2}\Bigl(1 - \alpha_{12}\alpha_{21}\Bigr)
\end{align*}

\schritt{Das Ergebnis}
\begin{center}
\begin{minipage}{0.9\textwidth}
\textbf{$D > 0$ genau dann, wenn $\alpha_{12}\alpha_{21} < 1$.} Und da $T$ ohnehin negativ ist, ist der Koexistenzpunkt genau unter dieser Bedingung stabil.
\end{minipage}
\end{center}

\begin{center}
\begin{minipage}{0.88\textwidth}
\kursiv{Damit ist Satz 5.3 sauber bewiesen, statt aus dem Bild abgelesen. Und man sieht zugleich, warum in Fall~4 ein Sattelpunkt vorliegt: Dort ist $\alpha_{12}\alpha_{21} > 1$, also $D < 0$.}
\end{minipage}
\end{center}
\end{merksatz}

\begin{merksatz}
\boxtitel{Mit den Zahlen aus Aufgabe 5.1}
$K_1=10$, $K_2=8$, $\alpha_{12}=0{,}6$, $\alpha_{21}=0{,}5$, $x^{*}=7{,}43$, $y^{*}=4{,}29$; dazu $r_1 = 1{,}0$, $r_2 = 0{,}8$.

\schritt{Die Matrix}
\[ J = \begin{pmatrix} -0{,}743 & -0{,}446\\ -0{,}215 & -0{,}429 \end{pmatrix} \]

\kursiv{Beispiel:} $-\frac{r_2\alpha_{21}y^{*}}{K_2} = -\frac{0{,}8 \cdot 0{,}5 \cdot 4{,}29}{8} = -\frac{1{,}716}{8} = -0{,}215$.

\schritt{Spur, Determinante, Diskriminante}
\[ T = -0{,}743 - 0{,}429 = -1{,}172 \]
\[ D = (-0{,}743)(-0{,}429) - (-0{,}446)(-0{,}215) = 0{,}3187 - 0{,}0959 = 0{,}2228 \]
\[ T^2 - 4D = 1{,}3736 - 0{,}8912 = 0{,}4824 > 0 \]

\schritt{Die Einordnung}
$D>0$, $T<0$, Diskriminante positiv: \textbf{stabiler Knoten.}

\schritt{Die Eigenwerte}
\[ \lambda_{1,2} = \frac{-1{,}172 \pm 0{,}6946}{2} \quad\Longrightarrow\quad \lambda_1 = -0{,}239,\quad \lambda_2 = -0{,}933 \]

\schritt{Die Probe}
$\lambda_1 + \lambda_2 = -1{,}172 = T$ \checkmark\\
$\lambda_1 \cdot \lambda_2 = 0{,}2230 \approx D$ \checkmark

\schritt{Die Deutung}
Die langsamere Annäherung erfolgt mit Rate $0{,}239$, also mit einer Halbwertszeit von
\[ \frac{\ln 2}{0{,}239} = 2{,}90 \ \text{Zeiteinheiten} . \]
Nach etwa vier solchen Zeiten --- rund zwölf Einheiten --- ist das Gleichgewicht praktisch erreicht. Das ist dieselbe Vier-Halbwertszeiten-Regel wie in Kapitel~1.
\end{merksatz}

\section{Der Bruchfall}

\begin{warnung}
\boxtitelr{Wo die Linearisierung nichts sagt}
\begin{center}
{\small
\renewcommand{\arraystretch}{1.52}
\begin{tabular}{@{}P{2.6cm}P{6.0cm}@{}}
\toprule
\textbf{Realteil null} & Der unentschiedene Fall aus Satz 6.2. Die Näherung sagt \enquote{Zentrum}, die Wirklichkeit kann eine langsam auslaufende oder aufklingende Spirale sein.\\
\textbf{Weit weg} & Die Aussage gilt nur \kursiv{nahe} dem Gleichgewicht. Wie nahe, sagt sie nicht.\\
\textbf{Mehrere Gleichgewichte} & Jedes wird einzeln untersucht; welches erreicht wird, bleibt offen.\\
\bottomrule
\end{tabular}}
\end{center}

\schritt{Zur zweiten Zeile}
Das ist die wichtigste Einschränkung. Ein stabiles Gleichgewicht kann ein winziges Einzugsgebiet haben.

\begin{center}
\begin{minipage}{0.88\textwidth}
\kursiv{Im Konkurrenzmodell Fall~4 sind beide Randzustände stabil --- und trotzdem landet man je nach Start bei verschiedenen. Die Stabilitätsuntersuchung sagt nichts darüber, wo die Grenze verläuft. Dafür braucht man die Bahn durch den Sattelpunkt, und die muss man ausrechnen.}
\end{minipage}
\end{center}
\end{warnung}

\begin{merksatz}
\boxtitel{Was sich beim Verändern eines Parameters ereignet}
Verschiebt man einen Parameter langsam, so wandert der Punkt $(T,D)$ durch das Bild aus Abschnitt 6.3. Überquert er eine Grenzlinie, ändert sich das Verhalten schlagartig. Man nennt das eine \Begriff{Bifurkation}.

\begin{center}
{\small
\renewcommand{\arraystretch}{1.5}
\begin{tabular}{@{}P{2.6cm}P{2.6cm}P{3.0cm}@{}}
\toprule
\textbf{Übergang} & \textbf{Grenze} & \textbf{Folge}\\
\midrule
Sattel-Knoten & $D = 0$ & zwei Gleichgewichte verschmelzen und verschwinden\\
Hopf & $T = 0$, $D>0$ & \textcolor{akzent}{aus Ruhe wird Schwingung}\\
Knoten-Strudel & $T^2 = 4D$ & nur die Form ändert sich\\
\bottomrule
\end{tabular}}
\end{center}

\begin{center}
\begin{minipage}{0.88\textwidth}
\kursiv{Die dritte Zeile ist harmlos --- die Stabilität ändert sich nicht. Die ersten beiden sind es nicht. Bei der Hopf-Bifurkation entsteht aus einem stabilen Ruhezustand eine anhaltende Schwingung; das ist der Weg, auf dem biologische Rhythmen in Modellen entstehen. Kapitel~12 zeigt es an der Erregung von Nervenzellen.}
\end{minipage}
\end{center}
\end{merksatz}

\begin{warnung}
\boxtitelr{Langsames Ankündigen}
Nähert sich ein System einer Bifurkation, so wird die Rückkehr nach einer Störung immer langsamer --- weil der maßgebliche Eigenwert gegen null geht.

\begin{center}
\begin{minipage}{0.9\textwidth}
\textbf{Das ist messbar, bevor der Übergang geschieht.} Wachsende Schwankungen und langsamere Erholung gelten als Frühwarnzeichen für einen bevorstehenden Umschlag.
\end{minipage}
\end{center}

\begin{center}
\begin{minipage}{0.88\textwidth}
\kursiv{Das Verfahren wird bei Ökosystemen, Klimakipppunkten und in der Medizin untersucht --- etwa bei der Frage, ob sich ein epileptischer Anfall oder ein Stimmungsumschwung ankündigt. Die Ergebnisse sind uneinheitlich: Das Zeichen tritt auf, ist aber wenig spezifisch. Kapitel~15 kommt darauf zurück.}
\end{minipage}
\end{center}
\end{warnung}

\section{Übungsaufgaben}

\begin{aufgabe}
\boxtitel{Aufgabe 6.1 --- Vier Systeme einordnen}
Bestimmen Sie Typ und Stabilität für die Matrizen
\[ J_1 = \begin{pmatrix} -2 & 1\\ 0 & -3\end{pmatrix},\quad
J_2 = \begin{pmatrix} 1 & 2\\ 3 & -1\end{pmatrix},\quad
J_3 = \begin{pmatrix} -1 & -4\\ 2 & -1\end{pmatrix},\quad
J_4 = \begin{pmatrix} 0 & -2\\ 3 & 0\end{pmatrix} . \]

\loesung
\schritt{$J_1$}
\[ T = -2 + (-3) = -5, \qquad D = (-2)(-3) - (1)(0) = 6 \]
\[ T^2-4D = 25 - 24 = 1 > 0 \]
$D>0$, $T<0$, Diskriminante positiv: \textbf{stabiler Knoten.}
\[ \lambda = \frac{-5 \pm 1}{2} = -2 \ \text{und}\ -3 \]
\kursiv{Probe:} Summe $-5$ \checkmark, Produkt $6$ \checkmark

\schritt{$J_2$}
\[ T = 1 + (-1) = 0, \qquad D = (1)(-1) - (2)(3) = -1 - 6 = -7 \]
$D < 0$: \textbf{Sattelpunkt, instabil.}
\[ T^2-4D = 0 + 28 = 28, \qquad \lambda = \frac{0 \pm 5{,}292}{2} = \pm2{,}646 \]
Ein positiver und ein negativer Eigenwert --- das Kennzeichen des Sattels.

\begin{center}
\begin{minipage}{0.86\textwidth}
\kursiv{Beachten Sie, dass hier $T = 0$ ist und der Punkt trotzdem eindeutig eingeordnet werden kann. Der unentschiedene Fall aus Satz 6.2 tritt nur ein, wenn $T=0$ \kursiv{und} $D>0$.}
\end{minipage}
\end{center}

\schritt{$J_3$}
\[ T = -1 + (-1) = -2, \qquad D = (-1)(-1) - (-4)(2) = 1 + 8 = 9 \]
\[ T^2-4D = 4 - 36 = -32 < 0 \]
$D>0$, $T<0$, Diskriminante negativ: \textbf{stabiler Strudel.}
\[ \lambda = \frac{-2 \pm \mathrm{i}\sqrt{32}}{2} = -1 \pm 2{,}828\,\mathrm{i} \]

\kursiv{Deutung:} Abklingrate $1$, Kreisfrequenz $2{,}828$, also Periode $2\pi/2{,}828 = 2{,}22$. Die Schwingung klingt mit der Halbwertszeit $\ln2/1 = 0{,}693$ ab --- deutlich schneller als eine Periode. Man sieht also kaum eine Schwingung, eher ein einmaliges Überschwingen.

\schritt{$J_4$}
\[ T = 0, \qquad D = 0 - (-2)(3) = 6 > 0 \]
$T=0$ und $D>0$: \textbf{Zentrum} --- und damit der \textcolor{akzent}{unentschiedene Fall}.
\[ \lambda = \pm\mathrm{i}\sqrt{6} = \pm2{,}449\,\mathrm{i} \]

\begin{center}
\begin{minipage}{0.86\textwidth}
\kursiv{Für die linearisierte Gleichung liegen geschlossene Bahnen vor. Ob das ursprüngliche nichtlineare System ebenfalls ein Zentrum hat, lässt sich daraus \kursiv{nicht} schließen.}
\end{minipage}
\end{center}

\schritt{Die Übersicht}
\begin{center}
\renewcommand{\arraystretch}{1.45}
\begin{tabular}{@{}P{1.2cm}P{1.4cm}P{1.4cm}P{1.6cm}P{2.4cm}@{}}
\toprule
& $T$ & $D$ & $T^2-4D$ & \textbf{Typ}\\
\midrule
$J_1$ & $-5$ & $6$ & $+1$ & stabiler Knoten\\
$J_2$ & $0$ & $-7$ & $+28$ & Sattel\\
$J_3$ & $-2$ & $9$ & $-32$ & stabiler Strudel\\
$J_4$ & $0$ & $6$ & $-24$ & Zentrum\\
\bottomrule
\end{tabular}
\end{center}
\end{aufgabe}

\begin{aufgabe}
\boxtitel{Aufgabe 6.2 --- Lotka-Volterra mit Kapazitätsgrenze}
Untersuchen Sie
\[ \frac{\dd x}{\dd t} = x\left(a\left(1-\frac{x}{K}\right) - by\right), \qquad \frac{\dd y}{\dd t} = y(-c+dx) \]
für $a=1{,}0$, $b=0{,}5$, $c=0{,}4$, $d=0{,}2$, $K=10$.

\loesung
\schritt{Schritt 1 --- Das Gleichgewicht}
Aus der zweiten Gleichung, mit $y \neq 0$:
\[ x^{*} = \frac{c}{d} = \frac{0{,}4}{0{,}2} = 2{,}0 \]
--- unverändert gegenüber Kapitel~4.

Aus der ersten, mit $x \neq 0$:
\[ a\left(1-\frac{x^{*}}{K}\right) = b\,y^{*} \quad\Longrightarrow\quad y^{*} = \frac{a}{b}\left(1-\frac{x^{*}}{K}\right) = \frac{1{,}0}{0{,}5}\left(1-\frac{2}{10}\right) = 2 \cdot 0{,}8 = 1{,}6 \]

\begin{center}
\begin{minipage}{0.86\textwidth}
\kursiv{Der Beutebestand bleibt bei $2{,}0$ --- die Überkreuz-Eigenschaft aus Kapitel~4 überlebt die Erweiterung. Der Räuberbestand sinkt von $2{,}0$ auf $1{,}6$, weil die Beute nun selbst begrenzt ist und weniger nachliefert.}
\end{minipage}
\end{center}

\schritt{Schritt 2 --- Die Ableitungen}
Ausmultipliziert: $f = ax - \frac{a}{K}x^2 - bxy$.
\[ \frac{\partial f}{\partial x} = a - \frac{2a}{K}x - by, \qquad \frac{\partial f}{\partial y} = -bx \]
\[ \frac{\partial g}{\partial x} = dy, \qquad \frac{\partial g}{\partial y} = -c+dx \]

\schritt{Schritt 3 --- Im Gleichgewicht einsetzen}
\[ \frac{\partial f}{\partial x} = 1{,}0 - \frac{2\cdot1{,}0}{10}\cdot2{,}0 - 0{,}5\cdot1{,}6 = 1{,}0 - 0{,}4 - 0{,}8 = -0{,}2 \]
\[ \frac{\partial f}{\partial y} = -0{,}5 \cdot 2{,}0 = -1{,}0 \]
\[ \frac{\partial g}{\partial x} = 0{,}2 \cdot 1{,}6 = 0{,}32 \]
\[ \frac{\partial g}{\partial y} = -0{,}4 + 0{,}2\cdot2{,}0 = 0 \]

\[ J = \begin{pmatrix} -0{,}2 & -1{,}0\\ 0{,}32 & 0 \end{pmatrix} \]

\schritt{Schritt 4 --- Einordnen}
\[ T = -0{,}2 + 0 = -0{,}2 \]
\[ D = (-0{,}2)(0) - (-1{,}0)(0{,}32) = 0{,}32 \]
\[ T^2-4D = 0{,}04 - 1{,}28 = -1{,}24 < 0 \]

$D>0$, $T<0$, Diskriminante negativ: \textbf{stabiler Strudel.}

\schritt{Schritt 5 --- Die Eigenwerte}
\[ \lambda = \frac{-0{,}2 \pm \mathrm{i}\sqrt{1{,}24}}{2} = -0{,}1 \pm 0{,}557\,\mathrm{i} \]

\schritt{Schritt 6 --- Die Deutung}
\begin{center}
\renewcommand{\arraystretch}{1.48}
\begin{tabular}{@{}P{2.8cm}P{2.4cm}P{2.4cm}@{}}
\toprule
& \textbf{Wert} & \\
\midrule
Abklingrate & $0{,}1\ \mathrm{h}^{-1}$ & \\
Halbwertszeit & $6{,}93$ h & $\ln2/0{,}1$\\
Kreisfrequenz & $0{,}557$ & \\
Periode & $11{,}3$ h & $2\pi/0{,}557$\\
\bottomrule
\end{tabular}
\end{center}

Die Schwingung hat eine Periode von etwa elf Stunden und klingt mit einer Halbwertszeit von sieben Stunden ab. Nach etwa vier Perioden ist sie praktisch verschwunden.

\schritt{Schritt 7 --- Der Vergleich mit Kapitel 4}
\begin{center}
\renewcommand{\arraystretch}{1.48}
\begin{tabular}{@{}P{2.8cm}P{2.4cm}P{2.4cm}@{}}
\toprule
& \textbf{ohne $K$} & \textbf{mit $K$}\\
\midrule
$T$ & $0$ & $\mathbf{-0{,}2}$\\
Typ & Zentrum & \textcolor{gruen}{stabiler Strudel}\\
Periode & $9{,}93$ h & $11{,}3$ h\\
\bottomrule
\end{tabular}
\end{center}

\begin{center}
\begin{minipage}{0.88\textwidth}
\kursiv{Damit ist der Bruchfall aus Kapitel~4 nachgewiesen statt behauptet: Die Kapazitätsgrenze verschiebt die Spur von null ins Negative, und aus dem Zentrum wird ein Strudel. \textbf{Die Ursache steckt allein im Term $-\frac{2a}{K}x$}, der bei $K \to \infty$ verschwindet --- dann wäre $T=0$ wie zuvor.}
\end{minipage}
\end{center}
\end{aufgabe}

\begin{aufgabe}
\boxtitel{Aufgabe 6.3 --- Eine Bifurkation aufspüren}
Für das System aus Aufgabe 6.2 sei $K$ veränderlich. Bei welchem $K$ ändert sich das Verhalten qualitativ, und was geschieht?

\loesung
\schritt{Schritt 1 --- Das Gleichgewicht allgemein}
\[ x^{*} = \frac{c}{d} = 2{,}0 \ \text{(von $K$ unabhängig)}, \qquad y^{*} = \frac{a}{b}\left(1-\frac{2}{K}\right) = 2\left(1-\frac{2}{K}\right) \]

\schritt{Schritt 2 --- Die erste Bedingung}
Damit $y^{*} > 0$ ist, muss
\[ 1 - \frac{2}{K} > 0 \quad\Longleftrightarrow\quad K > 2 . \]

\begin{center}
\begin{minipage}{0.9\textwidth}
\textbf{Bei $K = 2$ verschwindet der Räuber.} Ist die Kapazität der Beute nicht größer als der Bedarf des Räubers, so kann er nicht bestehen.
\end{minipage}
\end{center}

\schritt{Schritt 3 --- Die Spur allgemein}
\[ \frac{\partial f}{\partial x}\Big|_{*} = a - \frac{2a}{K}x^{*} - by^{*} \]
Mit $by^{*} = a\left(1-\frac{x^{*}}{K}\right)$:
\[ = a - \frac{2a x^{*}}{K} - a + \frac{ax^{*}}{K} = -\frac{a\,x^{*}}{K} \]

Da $\frac{\partial g}{\partial y} = 0$ im Gleichgewicht:
\[ T = -\frac{a\,x^{*}}{K} = -\frac{1{,}0 \cdot 2{,}0}{K} = -\frac{2}{K} \]

\schritt{Schritt 4 --- Die Determinante allgemein}
\[ D = 0 - (-b x^{*})(d y^{*}) = b\,d\,x^{*}y^{*} = 0{,}5 \cdot 0{,}2 \cdot 2{,}0 \cdot y^{*} = 0{,}2\,y^{*} \]
\[ = 0{,}2 \cdot 2\left(1-\frac{2}{K}\right) = 0{,}4\left(1-\frac{2}{K}\right) \]

\schritt{Schritt 5 --- Die Werte im Überblick}
\begin{center}
{\small
\renewcommand{\arraystretch}{1.48}
\begin{tabular}{@{}P{1.4cm}P{1.6cm}P{1.8cm}P{1.6cm}P{2.0cm}@{}}
\toprule
$K$ & $y^{*}$ & $T$ & $D$ & \textbf{Typ}\\
\midrule
$2{,}5$ & $0{,}40$ & $-0{,}800$ & $0{,}080$ & Knoten\\
$3{,}0$ & $0{,}67$ & $-0{,}667$ & $0{,}133$ & Knoten\\
$4{,}0$ & $1{,}00$ & $-0{,}500$ & $0{,}200$ & \textcolor{akzent}{Grenze}\\
$6{,}0$ & $1{,}33$ & $-0{,}333$ & $0{,}267$ & Strudel\\
$10{,}0$ & $1{,}60$ & $-0{,}200$ & $0{,}320$ & Strudel\\
$\infty$ & $2{,}00$ & $\mathbf{0}$ & $0{,}400$ & \textcolor{tief}{Zentrum}\\
\bottomrule
\end{tabular}}
\end{center}

\schritt{Schritt 6 --- Die Grenze bestimmen}
Der Übergang Knoten/Strudel liegt bei $T^2 = 4D$:
\[ \frac{4}{K^2} = 4 \cdot 0{,}4\left(1-\frac{2}{K}\right) = 1{,}6 - \frac{3{,}2}{K} \]
Mit $K^2$ multiplizieren:
\[ 4 = 1{,}6K^2 - 3{,}2K \quad\Longrightarrow\quad 1{,}6K^2 - 3{,}2K - 4 = 0 \]
\[ K^2 - 2K - 2{,}5 = 0 \quad\Longrightarrow\quad K = 1 \pm \sqrt{1+2{,}5} = 1 \pm 1{,}871 \]
Die positive Lösung: $K = 2{,}871$.

\kursiv{Kontrolle:} Bei $K=4$ ist $T^2 = 0{,}25$ und $4D = 0{,}80$ --- also bereits Strudel. Bei $K=2{,}5$ ist $T^2=0{,}64$ und $4D=0{,}32$ --- Knoten. Die Grenze liegt dazwischen. \checkmark

\schritt{Schritt 7 --- Die Zusammenfassung}
\begin{center}
{\small
\renewcommand{\arraystretch}{1.48}
\begin{tabular}{@{}P{2.6cm}P{5.8cm}@{}}
\toprule
$K \le 2$ & Der Räuber stirbt aus; die Beute strebt gegen $K$.\\
$2 < K < 2{,}87$ & Stabiler Knoten --- Annäherung ohne Schwingen.\\
$K > 2{,}87$ & Stabiler Strudel --- gedämpfte Schwingungen.\\
$K \to \infty$ & Zentrum --- ungedämpfte Schwingungen.\\
\bottomrule
\end{tabular}}
\end{center}

\schritt{Schritt 8 --- Die Deutung}
\begin{center}
\begin{minipage}{0.9\textwidth}
\textbf{Je reicher die Umwelt --- je größer $K$ ---, desto schwächer gedämpft sind die Schwingungen.}
\end{minipage}
\end{center}

\begin{center}
\begin{minipage}{0.88\textwidth}
\kursiv{Dieses Ergebnis ist der Kern des sogenannten \Begriff{Anreicherungsparadoxons}: Verbessert man die Lebensbedingungen der Beute, so wird das System \kursiv{instabiler}, nicht stabiler. Bei realistischeren Fraßfunktionen kippt es sogar ganz um --- die Schwingungen wachsen an, bis eine Art ausstirbt. Rosenzweig beschrieb das 1971; es gilt als eines der überraschendsten Ergebnisse der theoretischen Ökologie.}
\end{minipage}
\end{center}
\end{aufgabe}

\begin{geschichte}
\boxtitelg{Zwei Wege zur selben Frage}
Die Frage, ob ein Gleichgewicht stabil ist, stellte sich zuerst in der Himmelsmechanik: Bleibt das Sonnensystem, wie es ist?

\Person{Henri Poincaré} (1854--1912) entwickelte ab 1881 die Methoden, die dieses Kapitel benutzt --- Phasenporträt, Nullklinen, Einordnung der Gleichgewichte. Er tat es, weil er die Bahnen nicht ausrechnen konnte und deshalb einen Weg suchte, ihr Verhalten \kursiv{ohne} Lösung zu bestimmen.

\begin{center}
\begin{minipage}{0.9\textwidth}
\kursiv{Poincaré fand dabei auch die empfindliche Abhängigkeit von den Anfangsbedingungen --- siebzig Jahre vor Lorenz. In der Preisschrift von 1890 über das Dreikörperproblem beschreibt er Bahnen von einer Verwicklung, die er selbst nicht zeichnen zu können erklärte. Die Stelle wurde jahrzehntelang übersehen.}
\end{minipage}
\end{center}

\Person{Alexander Ljapunow} (1857--1918) bewies 1892 in seiner Dissertation, dass die Linearisierung tatsächlich entscheidet --- sofern kein Realteil null ist. Das ist genau der Vorbehalt aus Satz 6.2, und er stammt von ihm.

\begin{center}
\begin{minipage}{0.9\textwidth}
\kursiv{Ljapunow entwickelte daneben eine zweite Methode, die ohne Linearisierung auskommt: Man sucht eine Funktion, die entlang jeder Bahn abnimmt --- ähnlich der Erhaltungsgröße aus Satz 4.3, nur mit Ungleichung statt Gleichung. Diese \enquote{direkte Methode} gilt bis heute als das mächtigste Stabilitätswerkzeug.}
\end{minipage}
\end{center}

\schritt{Der Weg in die Biologie}
Poincaré und Ljapunow arbeiteten an Himmelsmechanik und Regelungstechnik. Dass ihre Werkzeuge auf Populationen passen, bemerkte man erst in den 1920er Jahren --- durch Lotka und Volterra.

\begin{center}
\begin{minipage}{0.9\textwidth}
\kursiv{Das ist ein wiederkehrendes Muster: Die mathematischen Mittel sind lange vor ihrer Anwendung da. Beim Chaos war es dieselbe Lage --- Poincaré 1890, Verwendung ab 1975.}
\end{minipage}
\end{center}
\end{geschichte}

\begin{forschung}
\boxtitelp{Stabilität heute}
\begin{center}
{\small
\renewcommand{\arraystretch}{1.45}
\begin{tabular}{@{}P{2.8cm}P{5.4cm}@{}}
\toprule
\textbf{Viele Arten} & Bei $n$ Arten ist $J$ eine $n\times n$-Matrix. Robert May zeigte 1972, dass zufällig verknüpfte große Systeme mit wachsender Artenzahl \kursiv{instabiler} werden --- entgegen der damaligen Lehrmeinung.\\
\textbf{Netzwerk\-struktur} & Die Auflösung des Widerspruchs: Echte Nahrungsnetze sind nicht zufällig verknüpft. Welche Strukturen stabilisieren, ist ein aktives Feld.\\
\textbf{Frühwarnzeichen} & Langsamere Erholung und wachsende Schwankungen vor Kipppunkten; geprüft an Seen, Klimadaten und medizinischen Zeitreihen.\\
\bottomrule
\end{tabular}}
\end{center}

Mays Ergebnis von 1972 ist ein Musterbeispiel dafür, wie ein mathematischer Befund ein Fach erschüttert. Die Ökologie hatte jahrzehntelang gelehrt, Vielfalt erzeuge Stabilität. May rechnete nach und fand das Gegenteil --- für zufällig verknüpfte Systeme.

\begin{center}
\begin{minipage}{0.88\textwidth}
\kursiv{Die Auflösung dauerte dreißig Jahre und lautet: Beides ist richtig, je nach Struktur. Vielfalt stabilisiert, wenn die Verknüpfungen bestimmten Mustern folgen --- viele schwache und wenige starke Wechselwirkungen, geschachtelte Nahrungsnetze. Zufällige Vielfalt destabilisiert. Das ist ein Ergebnis, das ohne die Matrixrechnung dieses Kapitels nicht formulierbar wäre.}
\end{minipage}
\end{center}
\end{forschung}

\begin{merksatz}
\boxtitel{Was Sie jetzt können}
\begin{center}
\renewcommand{\arraystretch}{1.5}
\begin{tabular}{@{}P{3.0cm}P{5.6cm}@{}}
\toprule
\textbf{Definition 6.1} & Jacobi-Matrix als Tangente in zwei Veränderlichen\\
\textbf{Satz 6.2} & Vorzeichen der Realteile entscheidet\\
\textbf{Satz 6.3} & stabil $\iff$ $T<0$ und $D>0$\\
\textbf{Die Landkarte} & Sattel, Knoten, Strudel, Zentrum\\
\textbf{Zwei Nachweise} & Lotka-Volterra ist ein Zentrum; Satz 5.3 bewiesen\\
\bottomrule
\end{tabular}
\end{center}

\begin{center}
\begin{minipage}{0.88\textwidth}
\kursiv{Damit endet Teil~II. Ab Kapitel~7 wenden wir dieselben Werkzeuge auf eine Frage an, bei der eine einzige Zahl über Ausbruch oder Abklingen entscheidet.}
\end{minipage}
\end{center}
\end{merksatz}
